% GENERATED FILE. Edit canonical lessons, then run npm run generate:books.
% canonical-source-hash: fnv1a64:49b90d2d13b7db18
% canonical-lessons: GE-C435-abkuerzung, GE-C435-abkuerzen, GE-C435-umweg, GE-C435-umkehren, GE-C435-verfahren

\chapter{Abkürzung oder Umweg}
\label{ch:ge-abkurzung-oder-umweg}
\hlchaptermodality{435}

\begin{chapteropening}
\textbf{By the end of this chapter:} \emph{I can talk about routes: the shortcut, taking a shortcut, the detour, turning back and getting lost by car.}

Complete the last lesson of chapter 435: I can talk about routes: the shortcut, taking a shortcut, the detour, turning back and getting lost by car.
\end{chapteropening}

\section[die Abkürzung]{die Abkürzung — the shortcut; the abbreviation}

\label{lesson:GE-C435-abkuerzung}



\begin{quote}
Before the new one: say the German for the turning, then the German for to branch off.
\end{quote}

\subsection*{You'll want to know: die Abkürzung}
\textbf{die Abkürzung} — \textquotedblleft{}the shortcut; the abbreviation\textquotedblright{}.

\begin{quote}
Durch den Park gibt es eine Abkürzung zum Bahnhof.
\end{quote}

\begin{quote}
\emph{There is a shortcut to the station through the park.}
\end{quote}

Plural: \emph{die Abkürzungen}. \emph{Usw.} is the \emph{Abkürzung} for \emph{und so weiter}.

\begin{cousinweb}
\emph{Abkürzung} is built on \emph{abkürzen}, from \emph{kurz} 'short', which goes back to Latin \emph{curtus} 'shortened' (EtymDB; Lewis \& Short).
\end{cousinweb}

\begin{grammarlens}[title={What you've built}]
One more word for telling, explaining and describing at B1.
\end{grammarlens}

\subsection*{Your turn}
\begin{itemize}
\raggedright
  \item \emph{Say it:} \emph{die Abkürzung}
  \item \emph{Say it:} \emph{die Abkürzung}, once more
  \item \emph{Say it:} a sentence of your own with \emph{die Abkürzung}
  \item \emph{Recall:} say the German for the turning, then the German for to branch off, then say \emph{die Abkürzung} again
\end{itemize}

\begin{tcolorbox}[breakable,colback=teal!4,colframe=teal!35!black,title={Before you move on}]
What does die Abkürzung mean? (\textquotedblleft{}The shortcut; the abbreviation\textquotedblright{}.) Say it once more.
\end{tcolorbox}
\section[abkürzen]{abkürzen — to take a shortcut; to shorten}

\label{lesson:GE-C435-abkuerzen}



\begin{quote}
Before the new one: say the German for to branch off, then the German for the shortcut.
\end{quote}

\subsection*{You'll want to know: abkürzen}
\textbf{abkürzen} — \textquotedblleft{}to take a shortcut; to shorten\textquotedblright{}.

\begin{quote}
Wir kürzen über den Parkplatz ab.
\end{quote}

\begin{quote}
\emph{We take a shortcut across the car park.}
\end{quote}

Separable: \emph{wir kürzen ab}, \emph{wir haben abgekürzt}.

\begin{cousinweb}
\emph{abkürzen} is built from \emph{ab} and \emph{kürzen} 'to shorten', from \emph{kurz} 'short', which goes back to Latin \emph{curtus} 'shortened' (EtymDB; Lewis \& Short).
\end{cousinweb}

\begin{grammarlens}[title={What you've built}]
One more word for telling, explaining and describing at B1.
\end{grammarlens}

\subsection*{Your turn}
\begin{itemize}
\raggedright
  \item \emph{Say it:} \emph{abkürzen}
  \item \emph{Say it:} \emph{abkürzen}, once more
  \item \emph{Say it:} a sentence of your own with \emph{abkürzen}
  \item \emph{Recall:} say the German for to branch off, then the German for the shortcut, then say \emph{abkürzen} again
\end{itemize}

\begin{tcolorbox}[breakable,colback=teal!4,colframe=teal!35!black,title={Before you move on}]
What does abkürzen mean? (\textquotedblleft{}To take a shortcut; to shorten\textquotedblright{}.) Say it once more.
\end{tcolorbox}
\section[der Umweg]{der Umweg — the detour, the long way round}

\label{lesson:GE-C435-umweg}



\begin{quote}
Before the new one: say the German for the shortcut, then the German for to take a shortcut.
\end{quote}

\subsection*{You'll want to know: der Umweg}
\textbf{der Umweg} — \textquotedblleft{}the detour, the long way round\textquotedblright{}.

\begin{quote}
Wegen der Baustelle machen wir einen Umweg.
\end{quote}

\begin{quote}
\emph{Because of the road works we take a detour.}
\end{quote}

Plural: \emph{die Umwege}. \emph{Auf Umwegen} — in a roundabout way.

\begin{cousinweb}
\emph{Umweg} is built from \emph{um} 'round, around' and \emph{Weg} 'way', inherited from Proto-Germanic \emph{wegaz} (EtymDB).
\end{cousinweb}

\begin{grammarlens}[title={What you've built}]
One more word for telling, explaining and describing at B1.
\end{grammarlens}

\subsection*{Your turn}
\begin{itemize}
\raggedright
  \item \emph{Say it:} \emph{der Umweg}
  \item \emph{Say it:} \emph{der Umweg}, once more
  \item \emph{Say it:} a sentence of your own with \emph{der Umweg}
  \item \emph{Recall:} say the German for the shortcut, then the German for to take a shortcut, then say \emph{der Umweg} again
\end{itemize}

\begin{tcolorbox}[breakable,colback=teal!4,colframe=teal!35!black,title={Before you move on}]
What does der Umweg mean? (\textquotedblleft{}The detour, the long way round\textquotedblright{}.) Say it once more.
\end{tcolorbox}
\section[umkehren]{umkehren — to turn back, to turn round}

\label{lesson:GE-C435-umkehren}



\begin{quote}
Before the new one: say the German for to take a shortcut, then the German for the detour.
\end{quote}

\subsection*{You'll want to know: umkehren}
\textbf{umkehren} — \textquotedblleft{}to turn back, to turn round\textquotedblright{}.

\begin{quote}
Die Straße ist gesperrt, deshalb kehren wir um.
\end{quote}

\begin{quote}
\emph{The road is closed, so we turn back.}
\end{quote}

Separable; Perfekt with \emph{sein}: \emph{wir sind umgekehrt}.

\begin{cousinweb}
\emph{umkehren} is inherited from Middle High German \emph{ummekēren} (EtymDB), built from \emph{um} 'round' and \emph{kehren} 'to turn'; Old English had a matching verb, \emph{ymbċierran}.
\end{cousinweb}

\begin{grammarlens}[title={What you've built}]
One more word for telling, explaining and describing at B1.
\end{grammarlens}

\subsection*{Your turn}
\begin{itemize}
\raggedright
  \item \emph{Say it:} \emph{umkehren}
  \item \emph{Say it:} \emph{umkehren}, once more
  \item \emph{Say it:} a sentence of your own with \emph{umkehren}
  \item \emph{Recall:} say the German for to take a shortcut, then the German for the detour, then say \emph{umkehren} again
\end{itemize}

\begin{tcolorbox}[breakable,colback=teal!4,colframe=teal!35!black,title={Before you move on}]
What does umkehren mean? (\textquotedblleft{}To turn back, to turn round\textquotedblright{}.) Say it once more.
\end{tcolorbox}
\section[verfahren]{verfahren — to get lost driving (mostly reflexive: sich verfahren); to proceed}

\label{lesson:GE-C435-verfahren}



\begin{quote}
Before the new one: say the German for the detour, then the German for to turn back.
\end{quote}

\subsection*{You'll want to know: verfahren}
\textbf{verfahren} — \textquotedblleft{}to get lost driving (mostly reflexive: sich verfahren); to proceed\textquotedblright{}.

\begin{quote}
Ohne Navi hat sich mein Vater verfahren.
\end{quote}

\begin{quote}
\emph{Without the satnav my father got lost.}
\end{quote}

Reflexive for getting lost: \emph{ich habe mich verfahren}. \emph{Das Verfahren} — the procedure.

\begin{cousinweb}
\emph{verfahren} is inherited from Middle High German \emph{vervarn} and Old High German \emph{firfaran} (EtymDB), built from \emph{ver-} and \emph{fahren} 'to travel'.
\end{cousinweb}

\begin{grammarlens}[title={What you've built}]
One more word for telling, explaining and describing at B1.
\end{grammarlens}

\subsection*{Your turn}
\begin{itemize}
\raggedright
  \item \emph{Say it:} \emph{verfahren}
  \item \emph{Say it:} \emph{verfahren}, once more
  \item \emph{Say it:} a sentence of your own with \emph{verfahren}
  \item \emph{Recall:} say the German for the detour, then the German for to turn back, then say \emph{verfahren} again
\end{itemize}

\begin{tcolorbox}[breakable,colback=teal!4,colframe=teal!35!black,title={Before you move on}]
What does verfahren mean? (\textquotedblleft{}To get lost driving (mostly reflexive: sich verfahren); to proceed\textquotedblright{}.) Say it once more.
\end{tcolorbox}
